% !TeX root = ../../Book1.tex
% 纲要标题：## 第1章　代数对象、映射与基本构造
% 纲要位置：计划纲要.md，第 131 行。

\chapter{代数对象、映射与基本构造}
\label{b1:ch:01}
\BookChapterNumber{b1:ch:01}

\begin{chapterabstract}
本章从熟悉的加法、复合与乘法出发，说明运算公理如何决定可以进行的推理。自由字、商集和生成子幺半群分别回答三个问题：如何记录运算表达式，如何把应当相同的表达式合并，以及如何用少量元素描述一个结构。映射的存在性与唯一性把这些构造联系起来，为下一章的群和商群准备共同语言。
\end{chapterabstract}

\begin{learninggoals}
\begin{itemize}
\item 区分结合律、交换律、单位元、逆元和消去律的作用，能用具体反例说明它们并不相互替代。
\item 构造自由幺半群和商幺半群，逐项核验运算的良定义及同态的延拓。
\item 分清生成元决定映射的唯一性与生成关系限制映射的存在性，并用泛性质证明指定同构的唯一性。
\end{itemize}
\end{learninggoals}

本章只要求集合、映射、整数和基本矩阵运算。读到新结构时，先找出它的元素、运算和相等标准，再问哪些映射保持这些资料。特别应亲自核对字的拼接与商集的计算；这些看似朴素的步骤将在商群、商环和模的构造中反复出现。

% !TeX root = ../../Book1.tex
% 纲要标题：### 1.1 运算与结构
% 纲要位置：计划纲要.md，第 133 行。

\section{运算与结构}
\label{b1:sec:0101}

% TODO: 按以下纲要要点撰写本节。

% 纲要内容（仅作写作计划，尚非教材正文）：
%
% * 集合上的一元、二元运算及其公理
% * 从整数加法、置换复合到矩阵乘法
% * 结合律、交换律、单位元与逆元的不同作用
%

\subsection{一元、二元运算及其公理}
\label{b1:subsec:0101:01}

待撰写。

\subsection{整数、置换与矩阵上的运算}
\label{b1:subsec:0101:02}

待撰写。

\subsection{结合律、交换律、单位元与逆元}
\label{b1:subsec:0101:03}

待撰写。

% !TeX root = ../../Book1.tex
% 纲要标题：### 1.2 半群与幺半群
% 纲要位置：计划纲要.md，第 139 行。

\section{半群与幺半群}
\label{b1:sec:0102}

% 纲要内容（仅作写作计划，尚非教材正文）：
%
% * 半群、幺半群及同态
% * 自由幺半群与字的拼接
% * 变换幺半群、消去律与群化的动机
%

\subsection{半群、幺半群及同态}
\label{b1:subsec:0102:01}

把结合律作为共同条件，便能同时讨论字的拼接、映射复合和矩阵乘法。由于这些运算通常不交换，本节使用乘法记号 \(ab\)，但不把因子顺序视为可以任意更换。

\begin{definition}\label{b1:def:monoid}
带有结合的二元运算的非空集合称为\term[banqun]{半群}[semigroup]。具有单位元的半群称为\term[yaobanqun]{幺半群}[monoid]，其单位元一般记为 \(1\)。幺半群 \(M,N\) 之间的\term[tongtai]{同态}[homomorphism]是满足
\[
f(ab)=f(a)f(b),\qquad f(1_M)=1_N
\]
的映射 \(f:M\rightarrow N\)。半群同态只要求乘法等式。
\end{definition}

单位保持条件不能从乘法保持条件自动得到。例如从单元素幺半群到整数乘法幺半群、把唯一元素送到 \(0\) 的映射保持乘法，却不保持单位。今后“幺半群同态”总包含保持单位的要求；这使空乘积与非空乘积具有相同的映射规则。

\ProseParagraphBreak
\((\mathbb N,+)\) 是以 \(0\) 为单位元的交换幺半群；正整数加法构成半群，却没有单位元。若 \(M\) 是幺半群，递归定义 \(a^0=1\)、\(a^{n+1}=a^na\)。结合律给出 \(a^{m+n}=a^ma^n\)：固定 \(m\) 对 \(n\) 归纳即可。于是任取 \(a\in M\)，映射 \(n\mapsto a^n\) 都是从 \((\mathbb N,+)\) 到 \(M\) 的同态。

\begin{proposition}\label{b1:prop:monoid-maps}
幺半群同态的复合仍是同态；恒等映射是同态。若同态 \(f:M\rightarrow N\) 是双射，则逆映射 \(f^{-1}\) 也保持乘法和单位。
\end{proposition}
\begin{proof}
对于同态 \(f:M\rightarrow N\)、\(g:N\rightarrow P\)，有
\((gf)(ab)=g(f(a)f(b))=(gf)(a)(gf)(b)\)，并且
\((gf)(1_M)=1_P\)。恒等映射的结论直接由定义得出。若 \(f\) 双射，写 \(x=f(a),y=f(b)\)，则
\(f^{-1}(xy)=f^{-1}(f(ab))=ab=f^{-1}(x)f^{-1}(y)\)，且
\(f^{-1}(1_N)=1_M\)。
\end{proof}

这样的双射同态称为\term[tonggou]{同构}[isomorphism]。它说明两个幺半群的乘法结构相同，只是元素的名称不同。若只比较底集合的大小，就会丢掉乘法信息；后面将反复用元素阶、交换性等性质区别同样大小的结构。

\subsection{自由幺半群与字的拼接}
\label{b1:subsec:0102:02}

给定集合 \(X\)，把它的元素看成字母。长度为 \(n\) 的\term[zi]{字}[word]是有序序列 \((x_1,\ldots,x_n)\)，通常写成 \(x_1\cdots x_n\)。长度零的字只有一个，称为空字，记为 \(\varepsilon\)。不同长度的字始终视为不同对象，即使 \(X\) 自身含有数或序列。

\begin{definition}\label{b1:def:free-monoid}
集合 \(X\) 上的\term[ziyou-yaobanqun]{自由幺半群}[free monoid]是所有有限字的集合
\[
X^*\coloneqq\coprod_{n\geq0}X^n,
\]
运算为前后拼接，单位元为 \(\varepsilon\)。记 \(\iota:X\rightarrow X^*\) 为把字母送到长度一的字的映射。
\end{definition}
\symidx{\(X^*\)：自由幺半群}

符号 \(\coprod\) 表示不交并，即保留每个字的长度标记。拼接后各位置的字母就是先列出第一个字，再列出第二个字。因此两次拼接的结果不依赖括号；空字不增加任何位置，确为单位元。例如 \(X=\{a,b\}\) 时，\((ab)(baa)=abbaa\)，而 \((baa)(ab)=baaab\)。所以当至少有两个字母时，自由幺半群并不交换。单字母的情形只有 \(\varepsilon,a,a^2,\ldots\)，按长度同构于 \((\mathbb N,+)\)。

\begin{theorem}[字母映射的延拓]\label{b1:thm:free-monoid}
任给幺半群 \(M\) 和集合映射 \(u:X\rightarrow M\)，存在唯一幺半群同态
\(\widehat u:X^*\rightarrow M\) 满足 \(\widehat u\iota=u\)。它由
\[
\widehat u(\varepsilon)=1_M,\qquad
\widehat u(x_1\cdots x_n)=u(x_1)\cdots u(x_n)
\]
给出。
\end{theorem}
\begin{proof}
一个非空字有唯一的字母序列，右端乘积由结合律不依赖括号，所以公式定义了映射。两个字拼接时，右端乘积恰是两个原乘积的乘积，故它保持乘法；空字的指定值保证保持单位。每个非空字都是长度一的字的乘积，任何满足要求的同态都只能取上述值；在空字上也只能取单位元，所以延拓唯一。
\end{proof}

例如给 \(a,b\) 指定两个同阶方阵 \(A,B\)，每个字便成为按既定次序进行的矩阵乘积。可能有不同字取同一矩阵，但在 \(X^*\) 中它们仍是不同的字。“自由”所排除的是额外规定的字之间的等式，并不要求每次代入后的取值都不同。

\subsection{变换幺半群、消去律与群化的动机}
\label{b1:subsec:0102:03}

集合 \(X\) 的所有自映射在复合下构成\term[bianhuan-yaobanqun]{变换幺半群}[transformation monoid]，记为 \(\operatorname{End}_{\mathrm{Set}}(X)\)，单位元为恒等映射。前节的移位映射说明，其中可有左逆而无右逆。由自映射到置换的限制，正是从一般操作转向可逆操作。

\begin{definition}\label{b1:def:cancellation}
幺半群满足\term[xiaoqulv]{消去律}[cancellation law]，是指
\(ax=ay\Rightarrow x=y\) 与 \(xa=ya\Rightarrow x=y\) 都成立。
\end{definition}

自由幺半群满足消去律，因为等字有相同的逐位字母，去掉相同前缀或后缀仍得等字。然而非空字没有逆元：拼接后的长度是两字长度之和，不可能成为零。因此消去律本身不等于可逆性。

\ProseParagraphBreak
对交换幺半群，可以像从自然数构造整数那样加入形式差。以下给出消去律充分发挥作用的一种情形。为了使公式清晰，暂以加法写 \(M\)，单位元记为 \(0\)。

\begin{proposition}[交换消去幺半群的群化]\label{b1:prop:commutative-completion}
设 \(M\) 是满足消去律的交换幺半群。在 \(M\times M\) 上规定
\[
(a,b)\sim(c,d)\quad\Longleftrightarrow\quad a+d=c+b.
\]
其等价类集合 \(K(M)\) 在
\([(a,b)]+[(c,d)]=[(a+c,b+d)]\) 下是交换幺半群，每个元素都有逆元 \([(b,a)]\)。映射
\(j(a)=[(a,0)]\) 是单射同态。
\end{proposition}
\begin{proof}
自反性和对称性直接成立。若 \(a+d=c+b\) 且 \(c+f=e+d\)，相加后消去 \(c+d\)，得到 \(a+f=e+b\)，故有传递性。若两对代表元分别等价，则把相应两条等式相加可知它们的逐项和等价，所以加法良定义。结合律与交换律来自 \(M\)，零元为 \([(0,0)]\)，而
\([(a+b,b+a)]=[(0,0)]\)，给出所述逆元。最后，\(j(a)=j(c)\) 当且仅当 \(a=c\)，并且 \(j\) 保持加法与零元。
\end{proof}

这里“等价类”的一般性质将在下一节系统说明；上述证明已经逐项给出了本构造需要的检验。对 \(M=\mathbb N\)，映射 \([(a,b)]\mapsto a-b\) 是到 \(\mathbb Z\) 的双射并保持加法。一般情形也有同样的映射原则：若 \(u:M\rightarrow A\) 是到一个每个元素都可逆的交换幺半群的同态，则
\([(a,b)]\mapsto u(a)-u(b)\) 是唯一延拓 \(u\) 的同态。良定义来自
\(a+d=c+b\) 蕴含 \(u(a)+u(d)=u(c)+u(b)\)；唯一性来自
\([(a,b)]=j(a)-j(b)\)。这说明新加的形式差除了使元素可逆，没有强加不必要的关系。

\ProseParagraphBreak
交换性在这里允许整理求和，消去律则用于等价关系的传递性和嵌入判断。非交换情形不能直接把同一公式中的加号换成乘号；一般半群嵌入群的问题有额外障碍，本节不作这种推广。

% !TeX root = ../../Book1.tex
% 纲要标题：### 1.3 等价关系与商集
% 纲要位置：计划纲要.md，第 145 行。

\section{等价关系与商集}
\label{b1:sec:0103}

% TODO: 按以下纲要要点撰写本节。

% 纲要内容（仅作写作计划，尚非教材正文）：
%
% * 等价关系、划分与商映射
% * 运算何时能够下降到商集
% * 同余关系与良定义性的逐项检验
%

\subsection{等价关系、划分与商映射}
\label{b1:subsec:0103:01}

待撰写。

\subsection{商集上的诱导运算}
\label{b1:subsec:0103:02}

待撰写。

\subsection{同余关系与良定义性}
\label{b1:subsec:0103:03}

待撰写。

% !TeX root = ../../Book1.tex
% 纲要标题：### 1.4 生成与泛性质
% 纲要位置：计划纲要.md，第 151 行。

\section{生成与泛性质}
\label{b1:sec:0104}

% 纲要内容（仅作写作计划，尚非教材正文）：
%
% * 生成子对象与生成元
% * 自由对象的映射延拓性质
% * 存在性、唯一性与典范同构的区别
%
% ---
%

\subsection{生成子对象与生成元}
\label{b1:subsec:0104:01}

要描述一个很大的代数对象，直接列出所有元素通常并不可行。更有效的办法是先挑选少量元素，再说明从这些元素可以进行哪些运算。生成问题因而分成两个层次：哪些元素能够得到，以及同一个元素有哪些不同的表达式。

\begin{definition}\label{b1:def:generated-monoid}
幺半群 \(M\) 的子集 \(N\) 若包含 \(1_M\) 且对乘法封闭，称为\term[zi-yaobanqun]{子幺半群}[submonoid]。给定 \(S\subseteq M\)，包含 \(S\) 的最小子幺半群称为 \(S\) \term[shengcheng]{生成}[generate]的子幺半群，记为 \(\langle S\rangle_{\mathrm{mon}}\)。若它等于 \(M\)，则称 \(S\) 为一个生成元集合。
\end{definition}
\symidx{\(\langle S\rangle_{\mathrm{mon}}\)：生成子幺半群}

\begin{proposition}\label{b1:prop:generated-monoid}
\(\langle S\rangle_{\mathrm{mon}}\) 存在，既等于所有包含 \(S\) 的子幺半群的交，又等于 \(S\) 中元素的所有有限乘积组成的集合，其中包括空乘积 \(1_M\)。
\end{proposition}
\begin{proof}
包含 \(S\) 的子幺半群至少有 \(M\) 自身。它们的交含有单位元；若两个元素属于交，它们属于每个子幺半群，因而乘积仍属于每一个，也属于交。这给出最小子幺半群。另一方面，所有有限乘积的集合包含 \(S\) 和单位元，并且两个有限乘积的乘积仍为有限乘积；任意包含 \(S\) 的子幺半群都必须包含这些乘积。故两种描述相同。
\end{proof}

例如在 \((\mathbb N,+)\) 中，\(\{2,3\}\) 生成
\(\{0,2,3,4,5,\ldots\}\)。因为 \(1\) 不能写成非负整数系数的 \(2,3\) 之和，它不在其中；每个偶数至少为 \(2\) 时是若干个 \(2\) 的和，每个奇数至少为 \(3\) 时是 \(3\) 加一个偶数。这里 \(6=2+2+2=3+3\) 显示生成表达式一般不唯一。

\begin{proposition}\label{b1:prop:maps-on-generators}
若 \(S\) 生成幺半群 \(M\)，则两个同态 \(f,g:M\rightarrow N\) 在 \(S\) 上相等便在 \(M\) 上相等。
\end{proposition}
\begin{proof}
对非空乘积 \(s_1\cdots s_r\)，同态的值分别为
\(f(s_1)\cdots f(s_r)\) 与 \(g(s_1)\cdots g(s_r)\)，逐因子相等；空乘积的像都是 \(1_N\)。这些元素已穷尽 \(M\)。
\end{proof}

这个命题只保证唯一性，不保证任意指定生成元的像都能形成同态。例如上例中，若想把 \(2\) 送到 \(1\in\mathbb N\)、把 \(3\) 送到 \(1\in\mathbb N\)，则 \(6\) 按两种表达式应同时送到 \(3\) 与 \(2\)，产生矛盾。生成关系正是延拓必须遵守的限制。

\subsection{自由对象的映射延拓性质}
\label{b1:subsec:0104:02}

\begin{definition}\label{b1:def:universal-property}
用对所有目标对象的映射存在性和唯一性来刻画一个构造的条件，称为该构造的\term[fan-xingzhi]{泛性质}[universal property]。例如 \(X^*\) 连同字母映射 \(\iota:X\rightarrow X^*\) 的泛性质是：到任意幺半群的字母映射都能唯一延拓为同态。
\end{definition}

\cref{b1:thm:free-monoid}已经证明这一性质。它同时处理所有目标幺半群，比列举若干可代入的数或矩阵更强。自由对象包含的资料是“对象加指定生成映射”，不是仅仅一个没有标记的幺半群。

\begin{proposition}[生成元与关系]\label{b1:prop:monoid-generators-relations}
设 \(R\) 是 \(X^*\) 中若干对字组成的集合。存在包含这些字对的最小同余关系 \(\equiv_R\)。商幺半群 \(X^*/{\equiv_R}\) 具有如下性质：它到 \(M\) 的同态，恰好对应那些使每对 \((v,w)\in R\) 都满足
\(\widehat u(v)=\widehat u(w)\) 的映射 \(u:X\rightarrow M\)。
\end{proposition}
\begin{proof}
把所有包含 \(R\) 的同余关系相交。这样的关系至少有把所有字视为等价的关系；交仍是等价关系，并且仍与左右拼接相容，故给出最小同余。若 \(u\) 满足所有指定等式，则 \(v\sim w\Longleftrightarrow\widehat u(v)=\widehat u(w)\) 是包含 \(R\) 的同余，因而包含 \(\equiv_R\)。由商幺半群的映射性质，\(\widehat u\) 唯一下降。反过来，商上的同态与商映射复合，在每对指定关系上必取相同值；限制到字母上恢复 \(u\)。这两个过程互为逆过程。
\end{proof}

例如在 \(\{a\}^*\) 中施加 \(a^2=a\)，任何正长度的字都等价于 \(a\)，所以商最多有 \(\varepsilon,a\) 两个元素。它们不会进一步相等：把 \(a\) 送到两元素幺半群 \(\{1,e\}\) 中满足 \(e^2=e\ne1\) 的 \(e\)，可区分二者。于是商正好是这一个两元素幺半群。向任意 \(M\) 指定它的同态，相当于在 \(M\) 中挑选一个满足 \(m^2=m\) 的元素。

\subsection{存在性、唯一性与典范同构的区别}
\label{b1:subsec:0104:03}

构造问题中，“存在”“唯一”和“典范”常被放在一句话里，但它们承担不同职责。存在性要求实际给出对象或论证其可取得；唯一性必须说明相对于哪些固定资料；典范性则强调映射由这些资料强制决定，不依赖额外选择。

\begin{theorem}[泛性质的唯一性]\label{b1:thm:universal-uniqueness}
若幺半群 \(F,F'\) 连同映射 \(i:X\rightarrow F\)、\(i':X\rightarrow F'\) 都满足自由幺半群的映射延拓性质，则存在唯一同构 \(\alpha:F\rightarrow F'\) 满足 \(\alpha i=i'\)。
\end{theorem}
\begin{proof}
由 \(F\) 的泛性质，\(i'\) 唯一延拓为同态 \(\alpha:F\rightarrow F'\)；由 \(F'\) 的泛性质，\(i\) 唯一延拓为同态 \(\beta:F'\rightarrow F\)。有
\((\beta\alpha)i=i\)，而恒等同态也在 \(i\) 上取值为 \(i\)。唯一性因此给出
\(\beta\alpha=\operatorname{id}_F\)。交换两个对象得到
\(\alpha\beta=\operatorname{id}_{F'}\)，所以 \(\alpha\) 是同构。任何满足指定等式的同构首先是延拓 \(i'\) 的同态，故只能是 \(\alpha\)。
\end{proof}

这个唯一同构称为由给定资料确定的\term[dianfan-tonggou]{典范同构}[canonical isomorphism]。它不意味着 \(F\) 与 \(F'\) 是同一个集合，也不意味着任意同构都唯一。例如两个字母的自由幺半群有交换两字母的自同构；只有要求分别保持指定字母时，允许的自同构才只剩恒等映射。

\ProseParagraphBreak
同样的双向构造方法适用于商幺半群和交换消去幺半群的群化：先用一方的泛性质构造映射，再用另一方反向构造，最后由唯一性证明复合为恒等。与猜测两个对象的元素对应相比，这种方法往往更短，也更能说明同构为何自然。使用时仍须先证明有关泛性质；把“显然具有泛性质”当作构造的代替，并不能解决良定义或存在性问题。


\begin{chaptersummary}
结合律使有限乘积可以省略括号，单位元处理空乘积，逆元则允许撤销运算。半群与幺半群把这些性质从具体算术中抽离出来，同时保留足够多的例子说明它们的差别。

生成提供表达式，自由对象允许任意指定生成元的像，同余关系规定哪些表达式必须取同一值。处理一个新商构造时，应先检查等价关系，再检查运算下降，最后检查所需映射的良定义、存在性与唯一性。下一章把每个元素均可逆这一条件加入幺半群，得到群。
\end{chaptersummary}

伴读时可参看 Dummit 与 Foote 的《Abstract Algebra》\cite{DummitFoote2003}及 Lang 的《Algebra》\cite{Lang2002}中关于代数结构与映射的讨论。以下习题练习构造新的运算、检验商映射，以及用反例辨别各项公理的作用。

\BookExercises

\begin{exercise}\label{b1:ex:01-affine}
固定实数 \(\alpha,\beta\)，在 \(\mathbb R\) 上定义 \(x*y=\alpha x+\beta y\)。求使运算满足结合律的全部 \((\alpha,\beta)\)，并在这些情形中判断何时有单位元、何时满足交换律。
\end{exercise}
\begin{solution}
展开得
\((x*y)*z=\alpha^2x+\alpha\beta y+\beta z\)，
\(x*(y*z)=\alpha x+\alpha\beta y+\beta^2z\)。对任意实数成立，当且仅当
\(\alpha^2=\alpha\)、\(\beta^2=\beta\)，故二者各为 \(0\) 或 \(1\)。若有双边单位元 \(e\)，从 \(e*y=y\) 比较 \(y\) 的系数得 \(\beta=1\)，从 \(x*e=x\) 得 \(\alpha=1\)，继而 \(e=0\)。交换律要求 \(\alpha=\beta\)，所以结合情形中只有 \((0,0)\) 与 \((1,1)\) 满足。前者处处取零而没有单位元，说明交换与结合仍不足以保证存在单位。
\end{solution}

\begin{exercise}\label{b1:ex:01-odd}
证明奇整数在乘法下构成幺半群，求其中所有有逆的元素。证明其消去律成立，但不能仅靠这些逆元生成整个幺半群。
\end{exercise}
\begin{solution}
奇数的乘积为奇数，\(1\) 是奇数，结合律来自整数乘法。若奇整数 \(a,b\) 满足 \(ab=1\)，则 \(a=b=1\) 或 \(a=b=-1\)，故有逆元素只有 \(\pm1\)。每个奇数非零，所以 \(ax=ay\) 在整数中推出 \(x=y\)，右消去同样成立。但 \(\pm1\) 的任意有限乘积仍为 \(\pm1\)，不能得到 \(3\)。因此一个消去幺半群既可能含有非平凡的可逆元素，也可能包含许多不能由这些可逆元素生成的元素。
\end{solution}

\begin{exercise}\label{b1:ex:01-length-mod}
设非空字母集 \(X\) 与整数 \(m\geq2\) 给定，规定两字等价当且仅当长度模 \(m\) 相同。描述所得商幺半群，并求任一类的逆元。
\end{exercise}
\begin{solution}
长度模 \(m\) 相同是等价关系，因为整数同余具有自反、对称和传递性。又因长度在拼接时相加，两对字分别等价时，它们的拼接也等价，故这是\cref{b1:def:congruence}中的同余关系。由\cref{b1:prop:operation-descends}，拼接能下降为商上的乘法。所有余数 \(0,\ldots,m-1\) 都能由同一个固定字母重复相应次数取得，因此商恰有 \(m\) 个元素，运算为余数相加再取模 \(m\)。长度余数为 \(r\) 的类的逆元是余数为 \(m-r\) 的类，若 \(r=0\) 则为自身。商映射可能把长度正的字送到单位类；这不表示自由幺半群中的原字已有逆元，而是施加关系后产生了新的可逆性。
\end{solution}

\begin{exercise}\label{b1:ex:01-transform-cancel}
设 \(X\) 至少含两个元素，\(f:X\rightarrow X\)。证明：对所有自映射 \(u,v\)，\(fu=fv\Rightarrow u=v\) 当且仅当 \(f\) 单射；\(uf=vf\Rightarrow u=v\) 当且仅当 \(f\) 满射。
\end{exercise}
\begin{solution}
若 \(f\) 单射，\(f\bigl(u(x)\bigr)=f\bigl(v(x)\bigr)\) 逐点推出 \(u(x)=v(x)\)。若不单射，取不同的 \(a,b\) 而 \(f(a)=f(b)\)，以 \(a,b\) 为值的两个常映射给出反例。若 \(f\) 满射，则任意 \(y\in X\) 可写成 \(f(x)\)，从 \(uf=vf\) 得 \(u(y)=v(y)\)。若不满射，取 \(z\notin f(X)\) 及不同的 \(a,b\in X\)，令 \(u\) 恒取 \(a\)，而 \(v\) 只在 \(z\) 处取 \(b\)、其余处取 \(a\)。于是 \(uf=vf\) 却 \(u\ne v\)。左右次序不同，所需的单射、满射条件也不同。
\end{solution}

\begin{exercise}\label{b1:ex:01-free-commutative}
设 \(X\) 为任意集合。令 \(C(X)\) 为所有仅在有限多个点非零的函数 \(a:X\rightarrow\mathbb N\)，运算为逐点相加。证明它具有交换幺半群中的自由映射延拓性质，并解释它与 \(X^*\) 的关系。
\end{exercise}
\begin{solution}
两个有限非零集合的并仍有限，所以加法封闭，零函数为单位，结合律和交换律逐点成立。记 \(\delta_x\) 为在 \(x\) 处取 \(1\)、其余取零的函数。对交换幺半群 \(M\) 和映射 \(u:X\rightarrow M\)，定义
\(\widehat u(a)=\prod_{a(x)\ne0}u(x)^{a(x)}\)。有限性使乘积存在，交换性使取点次序无关；指数相加公式证明它是同态。每个 \(a\) 唯一写成有限和 \(\sum_x a(x)\delta_x\)，故延拓唯一。把字送到各字母出现次数的函数得到满同态 \(X^*\rightarrow C(X)\)；其相等关系就是忽略字母次序、保留字母个数。
\end{solution}

\begin{exercise}\label{b1:ex:01-truncated-addition}
固定 \(n\geq1\)，在 \(T_n=\{0,1,\ldots,n\}\) 上定义 \(a\star b=\min(n,a+b)\)。证明这是幺半群，说明它是 \((\mathbb N,+)\) 的商，并证明任何到所有元素都有逆元的幺半群的同态都取常值单位元。
\end{exercise}
\begin{solution}
等式 \(\min\bigl(n,\min(n,a+b)+c\bigr)=\min(n,a+b+c)\) 证明结合律，零为单位。映射 \(q(k)=\min(n,k)\) 满，且 \(q(k+\ell)=q(k)\star q(\ell)\)、\(q(0)=0\)，故为幺半群同态。其同值关系 \(k\sim\ell\Longleftrightarrow q(k)=q(\ell)\) 与加法相容；由\cref{b1:thm:monoid-quotient}，公式 \([k]\mapsto q(k)\) 给出商到 \(T_n\) 的同态。它由同值关系的定义为单射，由 \(q\) 满射而满射，所以确实是所需同构。设 \(f:T_n\rightarrow G\) 保持单位与乘法，且 \(G\) 中元素都有逆。因为 \(n\star a=n\)，有 \(f(n)f(a)=f(n)\)，左乘 \(f(n)^{-1}\) 得 \(f(a)=1\)。尤其 \(T_n\) 不能单射嵌入这样的 \(G\)。失败源于截断运算丢失了消去律。
\end{solution}

\begin{exercise}\label{b1:ex:01-finite-cancellation}
证明满足左消去律的有限幺半群中，每个元素都有双边逆。给出无限情形的反例。
\end{exercise}
\begin{solution}
固定 \(a\)。左乘映射 \(x\mapsto ax\) 由左消去律为单射，有限性使其满射，故存在 \(b\) 使 \(ab=1\)。又有 \(a(ba)=(ab)a=a=a1\)，再用左消去得 \(ba=1\)。这为每个 \(a\) 提供双边逆。无限幺半群 \((\mathbb N,+)\) 满足左右消去律，但 \(1\) 没有加法逆元；其平移映射 \(x\mapsto1+x\) 单射而不满射。有限性的作用恰是把单射升级为满射。
\end{solution}

\begin{exercise}\label{b1:ex:01-infinitely-generated-submonoid}
在 \((\mathbb N^2,+)\) 中考虑
\(M=\bigl\{(0,0)\bigr\}\cup\bigl\{\,(a,b)\mid a\geq1,\ b\geq0\,\bigr\}\)。证明 \(M\) 是子幺半群但不能由有限多个元素生成，尽管 \(\mathbb N^2\) 可以。
\end{exercise}
\begin{solution}
零元属于 \(M\)；两个非零元素的和第一坐标至少为 \(2\)，仍在 \(M\)，所以它是子幺半群。每个 \((1,b)\) 若写成若干非零元素的和，由第一坐标相加可知只能有一个加数，且这个加数就是 \((1,b)\)。因而任意生成元集合都必须包含无限多个不同的 \((1,b)\)。另一方面，\(\mathbb N^2\) 由 \((1,0),(0,1)\) 生成。可见有限生成性不必传给子幺半群。
\end{solution}
